%============================================================
% CH2 · P2 · PROBLEM 2 — COUNT NODES WITH THE HIGHEST SCORE (LeetCode 2049)
%============================================================
\subsection{Problem 2.2.2: Count Nodes With the Highest Score}
\problemlabel{2.2.2}{ssec:ch2-highest-score}

% example 1 tree (0-based labels), per-vertex styles g2p2-h<k>
\tikzset{
    g2p2-h0/.style=gv, g2p2-h1/.style=gv, g2p2-h2/.style=gv, g2p2-h3/.style=gv, g2p2-h4/.style=gv,
    g2p2-H/.pic={
        \node[g2p2-h0] (h0) at (0,0) {0};
        \node[g2p2-h2] (h2) at (-1.1,-1) {2}; \node[g2p2-h4] (h4) at (1.1,-1) {4};
        \node[g2p2-h1] (h1) at (-1.7,-2) {1}; \node[g2p2-h3] (h3) at (-0.5,-2) {3};
        \foreach \a/\b in {0/2, 0/4, 2/1, 2/3} \draw[gtree] (h\a) -- (h\b);
    }
}

\problemcard{2.2.2}{Count Nodes With the Highest Score}{LeetCode 2049}{https://leetcode.com/problems/count-nodes-with-the-highest-score/}{\exmedium\ \lc{Medium}\ (the ``part above'' $n - \mathrm{sz}[v]$)}{Subtree sizes by post-order aggregation; the pieces left by deleting $v$ are its child subtrees and the part above of size $n - \mathrm{sz}[v]$; multiply in 64 bits}

\subsubsection*{1. Problem statement}

A \textbf{binary} tree has $n$ nodes $0..n-1$, rooted at 0, given by \texttt{parents} with $\texttt{parents}[0] = -1$. The \textbf{score} of a node: delete the node and its incident edges; the tree falls into one or more non-empty pieces; the score is the product of the piece sizes. Return how many nodes have the highest score.

\textbf{Constraints:} $2 \le n \le 10^5$; \texttt{parents} describes a valid binary tree (each node has at most two children). Note: $\texttt{parents}[i]$ may be larger than $i$.

\begin{table}[H]
\centering\small
\begin{tabular}{@{}lcl@{}}
\toprule
\textbf{parents} & \textbf{Output} & \textbf{Scores of nodes $0, 1, 2, \dots$} \\
\midrule
$[-1,2,0,2,0]$ & \texttt{3} & $3, 4, 2, 4, 4$ \\
$[-1,2,0]$ & \texttt{2} & $2, 2, 1$ \\
\bottomrule
\end{tabular}
\end{table}

\begin{diagrambox}
\begin{tikzpicture}[x=1cm, y=0.95cm]
    \begin{scope}
        \tikzset{g2p2-h1/.style=gvdst, g2p2-h3/.style=gvdst, g2p2-h4/.style=gvdst}
        \pic {g2p2-H};
        \foreach \v/\s/\d in {0/3/right, 1/4/left, 2/2/left, 3/4/right, 4/4/right} \node[gdist, text=deepblue, \d=0pt of h\v] {score \s};
        \node[font=\scriptsize\bfseries] at (0,0.9) {example 1: three nodes score 4};
    \end{scope}
    \begin{scope}[xshift=5.2cm]
        \tikzset{g2p2-h2/.style=gvdead, g2p2-h1/.style=gvA, g2p2-h3/.style=gvB, g2p2-h0/.style=gvsel, g2p2-h4/.style=gvsel}
        \pic {g2p2-H};
        \node[font=\scriptsize\bfseries] at (0,0.9) {delete 2: pieces $1 \cdot 1 \cdot 2 = 2$};
    \end{scope}
    \begin{scope}[xshift=10.4cm]
        \tikzset{g2p2-h0/.style=gvdead, g2p2-h2/.style=gvA, g2p2-h1/.style=gvA, g2p2-h3/.style=gvA, g2p2-h4/.style=gvB}
        \pic {g2p2-H};
        \node[font=\scriptsize\bfseries] at (0,0.9) {delete the root: $3 \cdot 1 = 3$};
    \end{scope}
\end{tikzpicture}
\figcap{Example 1. Each colour is one piece. Deleting a leaf leaves one piece of $n - 1 = 4$ nodes, which here is the maximum.}
\end{diagrambox}

\subsubsection*{2. The graph model}

\begin{modelbox}[Graph Model --- Count Nodes With the Highest Score]
\renewcommand{\arraystretch}{1.35}
\begin{tabularx}{\linewidth}{@{}>{\bfseries}l X@{}}
Vertices & the $n$ nodes; root 0 \\
Edges & tree edges $\texttt{parents}[v] \to v$; at most two children per node; labels give \emph{no} order ($\texttt{parents}[1] = 2$ in example 1) \\
Weights & none; every node counts 1 \\
Question & $\mathrm{score}(v) = \prod_{c \text{ child of } v} \mathrm{sz}[c] \cdot (n - \mathrm{sz}[v])$, the last factor omitted for the root; count the $v$ attaining the maximum \\
Size & $n \le 10^5$: \bigO{n}; a path of $10^5$ nodes is legal, so the depth can be $10^5$; scores reach $3.7\cdot10^{13}$ \\
\end{tabularx}
\end{modelbox}

\begin{diagrambox}
\begin{tikzpicture}[x=1cm, y=1cm]
    \node[gvcur] (v) at (0,0) {$v$};
    \coordinate (p) at (0,1.0);
    \draw[gfaded, dashed] (p) -- (v);
    \draw[draw=softpurple, fill=optimizedbg, rounded corners=8pt] (-2.6,1.0) rectangle (2.6,2.4);
    \node[font=\scriptsize, text width=4.8cm, align=center] at (0,2.05) {everything outside the subtree of $v$: $n - \mathrm{sz}[v]$ nodes};
    \draw[draw=primaryblue, fill=primaryblue!12] (-1.2,-0.6) -- (-2.1,-2.0) -- (-0.3,-2.0) -- cycle;
    \draw[draw=recursionred, fill=dupfill] (1.2,-0.6) -- (0.3,-2.0) -- (2.1,-2.0) -- cycle;
    \draw[gfaded, dashed] (v) -- (-1.2,-0.6); \draw[gfaded, dashed] (v) -- (1.2,-0.6);
    \node[font=\scriptsize] at (-1.2,-1.55) {$\mathrm{sz}[\ell]$};
    \node[font=\scriptsize] at (1.2,-1.55) {$\mathrm{sz}[r]$};
    \node[note, text width=5.6cm, anchor=west, align=left] at (3.2,0) {$\mathrm{score}(v) = \mathrm{sz}[\ell] \cdot \mathrm{sz}[r] \cdot (n - \mathrm{sz}[v])$\\[2pt]missing child: factor omitted (not 0)\\root: no part above (factor omitted, not 0)\\all three sizes add up to $n - 1$};
\end{tikzpicture}
\figcap{The model: at most three pieces, all of them readable from the subtree sizes.}
\end{diagrambox}

\subsubsection*{3. Key observations}

\begin{itemize}[leftmargin=*]
    \item \textbf{The pieces are known from sizes alone.} The children's pieces are their subtrees; the remaining piece is everything else, $n - \mathrm{sz}[v]$ nodes. One post-order pass computes all $\mathrm{sz}$, then every score is \bigO{1}.
    \item \textbf{Empty pieces do not count.} ``One or more non-empty subtrees'': a missing child and the empty part above the root contribute nothing --- multiplying by $0$ would make the score 0.
    \item \textbf{64 bits are needed.} The three sizes add up to $n - 1$, so the product is at most $\left(\frac{n-1}{3}\right)^3 \approx 3.7\cdot10^{13}$ for $n = 10^5$ (AM--GM), above $2^{31} - 1 \approx 2.1\cdot10^9$ but far below $2^{63}$. The bound is attained: a chain of $33\,333$ nodes from the root ending at $m$, and two chains of $33\,333$ below $m$, give $\mathrm{score}(m) = 33\,333^3 = 37\,035\,925\,937\,037$. The smallest tree that overflows \texttt{int} has $n = 3872$.
    \item \textbf{The labels are no order.} $\texttt{parents}[i] > i$ is allowed, so a backward loop over the labels is wrong; build children lists and a BFS order.
\end{itemize}

\subsubsection*{4. Approaches}

\begin{bruteforcebox}
\textbf{Idea: really delete each node.} For every $v$, block it and run a BFS from every remaining unvisited node; multiply the component sizes.

\begin{lstlisting}[style=cpp]
#include <bits/stdc++.h>
using namespace std;
typedef long long ll;

class SolutionBrute {
public:
    int countHighestScoreNodes(vector<int>& parents) {
        int n = parents.size();
        vector<vector<int>> adj(n);
        for (int v = 1; v < n; v++) { adj[v].push_back(parents[v]); adj[parents[v]].push_back(v); }
        vector<ll> score(n, 1);
        for (int del = 0; del < n; del++) {
            vector<char> seen(n, 0);
            seen[del] = 1;                              // the deleted node blocks the search
            for (int s = 0; s < n; s++) {
                if (seen[s]) continue;
                ll piece = 0;
                vector<int> q(1, s);
                seen[s] = 1;
                for (size_t i = 0; i < q.size(); i++) {
                    piece++;
                    for (int w : adj[q[i]]) if (!seen[w]) { seen[w] = 1; q.push_back(w); }
                }
                score[del] *= piece;
            }
        }
        ll best = *max_element(score.begin(), score.end());
        return count(score.begin(), score.end(), best);
    }
};
\end{lstlisting}

\begin{complexitybox}
\textbf{Time:} $n$ searches of \bigO{n}: $\Theta(n^2) = 10^{10}$ for $n = 10^5$. \textbf{Space:} \bigO{n}. \textbf{Verdict:} \tle; it is the reference: it does not know anything about sizes or the part above.
\end{complexitybox}
\end{bruteforcebox}

\begin{algorithmbox}[Approach 2: BFS order, sizes backwards, scores from sizes]
\textbf{Idea.} Children lists; a BFS order from 0 (parents first); sizes by the backward loop (template (2) of the pattern); then $\mathrm{score}(v)$ = product of the children's sizes, times $n - \mathrm{sz}[v]$ unless $v = 0$.

\begin{lstlisting}[style=cpp]
#include <bits/stdc++.h>
using namespace std;
typedef long long ll;

class Solution {
public:
    int countHighestScoreNodes(vector<int>& parents) {
        int n = parents.size();
        vector<vector<int>> children(n);
        for (int v = 1; v < n; v++) children[parents[v]].push_back(v);
        vector<int> order(1, 0);                       // BFS order: parents before children
        for (size_t i = 0; i < order.size(); i++)
            for (int c : children[order[i]]) order.push_back(c);
        vector<ll> sz(n, 1);
        for (int i = n - 1; i >= 1; i--) sz[parents[order[i]]] += sz[order[i]];
        ll best = -1;
        int count = 0;
        for (int v = 0; v < n; v++) {
            ll score = 1;
            for (int c : children[v]) score *= sz[c];  // the parts below v
            if (v != 0) score *= n - sz[v];            // the part above v
            if (score > best) { best = score; count = 1; }
            else if (score == best) count++;
        }
        return count;
    }
};
\end{lstlisting}

\begin{complexitybox}
\textbf{Time:} \bigO{n}. \textbf{Space:} \bigO{n}, no recursion --- the path of $10^5$ nodes is handled (answer 2: the two ends score $n - 1$). \textbf{Verdict:} returns 3 and 2 on the two examples.
\end{complexitybox}
\end{algorithmbox}

\begin{invariantbox}[Why the scores are right]
\textbf{Sizes.} In BFS order every node appears after its parent, so in the backward loop every node is handled after all its children: $\mathrm{sz}[v]$ is final when $v$ hands it up (the pattern's invariant).

\textbf{Pieces.} Removing $v$ removes the edges to its children and to its parent. Each child $c$ keeps its whole subtree, which no longer touches anything else: one piece of $\mathrm{sz}[c]$ nodes. All other nodes --- the nodes outside the subtree of $v$ --- are still connected to each other, because the path between any two of them in the tree does not pass through $v$ (it would have to enter and leave the subtree of $v$ through the single edge above $v$). That is one piece of $n - \mathrm{sz}[v]$ nodes if this number is positive, i.e.\ if $v$ is not the root. These are all the pieces, so the product is the score. \hfill$\blacksquare$
\end{invariantbox}

\begin{alternativebox}[Approach 3: peel the leaves]
Another order with ``children first'': keep, for every node, the number of children whose size is still unknown; start a queue with the leaves; when a node is taken from the queue its size is final, it multiplies its parent's score by its size, and the parent joins the queue when its last child is done. This is Kahn's algorithm (Chapter~5) on the tree directed towards the root, and it needs no BFS from the root.

\begin{lstlisting}[style=cpp]
#include <bits/stdc++.h>
using namespace std;
typedef long long ll;

class SolutionPeel {
public:
    int countHighestScoreNodes(vector<int>& parents) {
        int n = parents.size();
        vector<int> pending(n, 0);                     // children whose size is not yet known
        for (int v = 1; v < n; v++) pending[parents[v]]++;
        vector<ll> sz(n, 1), score(n, 1);
        vector<int> q;
        for (int v = 0; v < n; v++) if (pending[v] == 0) q.push_back(v);
        for (size_t i = 0; i < q.size(); i++) {
            int v = q[i];                              // all children of v are finished
            if (v == 0) continue;
            score[v] *= n - sz[v];
            int p = parents[v];
            sz[p] += sz[v];
            score[p] *= sz[v];
            if (--pending[p] == 0) q.push_back(p);
        }
        ll best = *max_element(score.begin(), score.end());
        return count(score.begin(), score.end(), best);
    }
};
\end{lstlisting}
\begin{complexitybox}
\textbf{Time:} \bigO{n}. \textbf{Space:} \bigO{n}. \textbf{Verdict:} returns 3 and 2 on the examples; agrees with Approaches 1 and 2 on all random trees. The queue order on example 1 is $1, 3, 4, 2, 0$ --- a ``reverse BFS'' from the leaves.
\end{complexitybox}
\end{alternativebox}

\begin{tricksbox}[Approach 4: Implementation Refinements]
\begin{itemize}[leftmargin=*]
    \item \textbf{Recursive DFS.} \texttt{sz[v] = 1 + dfs(left) + dfs(right)} is the shortest code, but a path of $10^5$ nodes recurses $10^5$ deep; LeetCode's stack is not documented, so the BFS order is the safe choice.
    \item \textbf{Count while computing.} Keep \texttt{best} and \texttt{count} in one pass; ties are compared exactly because the scores are integers (never use \texttt{double} for products of this size).
    \item \textbf{Children lists are optional.} For a binary tree, two arrays \texttt{left}, \texttt{right} with $-1$ for ``none'' are enough.
\end{itemize}
\end{tricksbox}

\subsubsection*{5. Trace}

\begin{tracebox}[{Example 1: $\texttt{parents} = [-1,2,0,2,0]$, $n = 5$}]
\begin{center}
\begin{tikzpicture}[x=1.05cm]
    \node[font=\scriptsize\bfseries, anchor=east] at (-0.6,0) {order};
    \foreach \v [count=\i from 0] in {0,2,4,1,3} \node[cell, minimum width=9mm] (o\i) at (\i,0) {\v};
    \node[font=\scriptsize\bfseries, anchor=east] at (-0.6,-0.75) {parent};
    \foreach \p [count=\i from 0] in {--,0,0,2,2} \node[idx] at (\i,-0.75) {\p};
    \draw[->, thick, primaryblue] (4.3,-1.25) -- node[below, font=\scriptsize]{backward: $\mathrm{sz}[\mathrm{par}] \mathrel{+}= \mathrm{sz}[v]$} (-0.3,-1.25);
    \node[note, anchor=west, text width=3.8cm, align=left] at (5.0,-0.4) {BFS from 0 with \texttt{children[0] = 2 4}, \texttt{children[2] = 1 3}};
\end{tikzpicture}
\end{center}
\begin{center}
\begin{tikzpicture}[x=0.95cm, y=0.85cm]
    \tikzset{gv/.append style={minimum size=6.5mm, font=\scriptsize\bfseries}}
    \begin{scope}
        \tikzset{g2p2-h3/.style=gvdone, g2p2-h1/.style=gvdone, g2p2-h4/.style=gvdone}
        \pic {g2p2-H};
        \foreach \v/\s/\d in {0/1/right, 1/1/left, 2/1/left, 3/1/right, 4/1/right} \node[gdist, text=primaryblue, \d=0pt of h\v] {\s};
        \node[font=\scriptsize\bfseries] at (0,0.85) {1. leaves 3, 1, 4: final 1};
    \end{scope}
    \begin{scope}[xshift=4.6cm]
        \tikzset{g2p2-h3/.style=gvdone, g2p2-h1/.style=gvdone, g2p2-h4/.style=gvdone, g2p2-h2/.style=gvcur}
        \pic {g2p2-H};
        \foreach \v/\s/\d in {0/2/right, 1/1/left, 2/3/left, 3/1/right, 4/1/right} \node[gdist, text=primaryblue, \d=0pt of h\v] {\s};
        \node[font=\scriptsize\bfseries] at (0,0.85) {2. node 2: $1 + 1 + 1 = 3$};
    \end{scope}
    \begin{scope}[xshift=9.2cm]
        \tikzset{g2p2-h3/.style=gvdone, g2p2-h1/.style=gvdone, g2p2-h4/.style=gvdone, g2p2-h2/.style=gvdone, g2p2-h0/.style=gvdone}
        \pic {g2p2-H};
        \foreach \v/\s/\d in {0/5/right, 1/1/left, 2/3/left, 3/1/right, 4/1/right} \node[gdist, text=primaryblue, \d=0pt of h\v] {\s};
        \node[font=\scriptsize\bfseries] at (0,0.85) {3. root: $\mathrm{sz}[0] = 5$};
    \end{scope}
\end{tikzpicture}
\end{center}
\begin{center}\small
\begin{tabular}{@{}cccccl@{}}
\toprule
$v$ & children (sizes) & $\mathrm{sz}[v]$ & $n - \mathrm{sz}[v]$ & score & best / count \\
\midrule
0 & 2 (3), 4 (1) & 5 & --- (root) & $3 \cdot 1 = 3$ & 3 / 1 \\
1 & --- & 1 & 4 & $4$ & 4 / 1 \\
2 & 1 (1), 3 (1) & 3 & 2 & $1 \cdot 1 \cdot 2 = 2$ & 4 / 1 \\
3 & --- & 1 & 4 & $4$ & 4 / 2 \\
4 & --- & 1 & 4 & $4$ & 4 / 3 \\
\bottomrule
\end{tabular}
\end{center}
Return \textbf{3}. The backward loop hands up in the order $3, 1, 4, 2$; the sizes after each step were $\mathrm{sz} = (1,1,2,1,1)$, $(1,1,3,1,1)$, $(2,1,3,1,1)$, $(5,1,3,1,1)$. In snapshot 2 the root shows the partial size 2 (node 4 has arrived, node 2 not yet).
\end{tracebox}

\subsubsection*{6. Test suite}

\begin{testbox}
\centering\small
\begin{tabular}{@{}p{5.6cm}cp{6.0cm}@{}}
\toprule
\textbf{parents} & \textbf{Output} & \textbf{Why} \\
\midrule
$[-1,0]$ & \texttt{2} & both nodes score 1 \\
$[-1,0,0]$ & \texttt{2} & root $1 \cdot 1 = 1$, leaves $2$ \\
$[-1,0,1,2,3]$ (path) & \texttt{3} & scores $4,3,4,3,4$: ends and middle \\
$[-1,0,0,1,1,2,2]$ (complete) & \texttt{1} & root $3 \cdot 3 = 9$ beats leaves (6) \\
path, $n = 10^5$ & \texttt{2} & depth $10^5$: the two ends score $99\,999$ \\
three arms of $33\,333$, $n = 10^5$ & \texttt{1} & $\mathrm{score}(m) = 37\,035\,925\,937\,037$; with \texttt{int} the program returns \texttt{3} \\
\bottomrule
\end{tabular}
\end{testbox}
All outputs from Approach~2 (Approach~3 and the brute force agree on the small ones). In addition, the three approaches agreed on $5000$ random binary trees with $n \le 13$ and shuffled labels (so that parents can have larger labels than children).

\subsubsection*{7. Common mistakes}

\begin{warningbox}[Mistake 1 --- \texttt{int} scores]
Replacing \texttt{long long} by \texttt{int} still passes both examples and every small test --- and returns \texttt{3} instead of \texttt{1} on the three-arm tree, where the products wrap around. The first tree where it can happen has only $3872$ nodes.
\end{warningbox}

\begin{warningbox}[Mistake 2 --- multiplying by an empty piece]
\begin{lstlisting}[style=snippet]
ll score = sz[l] * sz[r] * (n - sz[v]);     // WRONG: an absent child or the root gives 0
ll score = 1;                               // RIGHT: multiply only the non-empty pieces
for (int c : children[v]) score *= sz[c];
if (v != 0) score *= n - sz[v];
\end{lstlisting}
\end{warningbox}

\begin{warningbox}[Mistake 3 --- looping over the labels]
\texttt{for (v = n-1; v >= 1; v--) sz[parents[v]] += sz[v];} is template (3) of the pattern, valid only when $\texttt{parents}[v] < v$. In example 1, node 1 hands its size to node 2 \emph{after} node 2 has already handed its own size to the root, so the root ends with $\mathrm{sz}[0] = 4$ instead of 5 (the answer survives by luck). On the path $[-1,2,3,0]$ ($0 - 3 - 2 - 1$) the wrong sizes change the answer: \texttt{1} instead of \texttt{2}. Among $5000$ random small binary trees with shuffled labels, the label loop gave a wrong answer on $881$.
\end{warningbox}

\subsubsection*{8. Takeaways}

\begin{itemize}[leftmargin=*]
    \item Deleting $v$ leaves the child subtrees and one piece of $n - \mathrm{sz}[v]$: one size pass answers every deletion.
    \item Bound products before choosing the type: three parts summing to $n - 1$ give at most $((n-1)/3)^3$.
    \item When the labels are not an order, make one: BFS from the root, or peel leaves.
\end{itemize}

\begin{variantsbox}
\variant{Smallest largest piece}{\varnone}
  {find the node whose deletion leaves the smallest largest piece (a centroid).}
  {the same pieces; maximum instead of product.}
  {template (4) of the pattern: $\max(\max_c \mathrm{sz}[c],\ n - \mathrm{sz}[v])$; on the tree $T$ it picks vertex 1 (largest piece 3).}
\variant{Any tree}{\varnone}
  {nodes may have any number of children.}
  {up to $\deg(v)$ pieces; the product can exceed 64 bits.}
  {the same loop; compare $\sum \log$ of the piece sizes (with a tolerance) instead of products. The log comparison picked the same count as exact products on $3000$ random trees with $n \le 40$.}
\variant{Number of separated pairs}{\varnone}
  {for every node, how many pairs of other nodes lose their connection when it is deleted.}
  {pairs taken from two different pieces.}
  {$\binom{n-1}{2} - \sum_{\text{pieces}} \binom{s}{2}$, all pieces from sizes; matched a brute force (delete, BFS, count pairs) on the same $3000$ random trees.}
\variant{Fuel through an edge}{\varin{https://leetcode.com/problems/minimum-fuel-cost-to-report-to-the-capital/}}
  {everybody drives to the root in cars of $s$ seats.}
  {the edge above $v$ carries $\mathrm{sz}[v]$ people.}
  {sum $\lceil \mathrm{sz}[v] / s \rceil$ over all edges (Exercise E2.3).}
\end{variantsbox}
